\documentclass[12pt, a4paper]{article}

%=============== PAQUETES ===============
\usepackage[utf8]{inputenc}
\usepackage[T1]{fontenc}
\usepackage[spanish,es-noshorthands]{babel}
\usepackage{geometry}
\usepackage{graphicx}
\usepackage{amsmath}
\usepackage{xcolor}
\usepackage{circuitikz}
\usepackage{enumitem}
\usepackage{float}
\usepackage{hyperref}
\usetikzlibrary{arrows.meta, babel}

%=============== CONFIGURACIÓN ===============
\geometry{a4paper, margin=2.5cm}

\hypersetup{
    colorlinks=true,
    linkcolor=blue,
    urlcolor=cyan,
    pdftitle={Circuito de Entrada - SMPS},
}

%=============== DOCUMENTO ===============
\begin{document}

\title{\textbf{Componentes del Circuito de Entrada de una Fuente de Alimentación Conmutada (SMPS)}}
\author{Apuntes de Electrónica}
\date{\today}
\maketitle

\section*{Introducción}
Los componentes de una fuente de alimentación conmutada (SMPS) desde la entrada de la red eléctrica hasta el rectificador primario de onda completa son los siguientes:

\begin{itemize}[itemsep=1.5ex]
    \item \textbf{Fusible:} Es el primer elemento de la entrada y protege al circuito abriéndose si la corriente supera su límite establecido. Debe ser de "fusión lenta" (anti transitorios) porque al conectar la fuente a la red se produce un pico de corriente alto al que debe sobrevivir. A veces, los fabricantes utilizan un \textbf{fusistor}, que es una resistencia especial de muy bajo valor óhmico diseñada para actuar como fusible.
    
    \item \textbf{Filtro EMC (Compatibilidad Electromagnética):} Su propósito es aislar el ruido eléctrico generado por interferencias electromagnéticas, tanto para que no entren desde la red hacia la fuente, como para que la fuente no emita ruido a la instalación eléctrica. Consta principalmente de:
    \begin{itemize}[itemsep=1ex, topsep=1ex]
        \item \textbf{Condensadores de seguridad:} Los de \textbf{tipo X} se conectan entre las dos líneas de alimentación y los de \textbf{tipo Y} se conectan entre la línea y la tierra. Su función es derivar las señales nocivas de alta frecuencia. Son de tipo "autorregenerable", lo que significa que en caso de un arco eléctrico destruyen solo una pequeña zona de su material en lugar de crear un cortocircuito peligroso.
        \item \textbf{Bobinas acopladas (Modo Común):} Normalmente se enrollan sobre un núcleo de ferrita. Se oponen con firmeza al paso de las señales indeseadas de alta frecuencia. Se enrollan deliberadamente en sentidos opuestos para que las tensiones continuas anulen sus campos magnéticos entre sí, evitando que el núcleo de ferrita se sature y pierda sus propiedades magnéticas.
    \end{itemize}

    \item \textbf{Resistencias de descarga:} Son resistencias ubicadas en la entrada (frecuentemente varias en serie) que garantizan la descarga eléctrica de los condensadores cuando se apaga o desconecta el equipo, previniendo posibles descargas al manipular la placa.

    \item \textbf{Interruptor de encendido o puente:} Los terminales de entrada de algunas placas pueden incluir un interruptor (ON/OFF) o bien un puente de alambre soldado para dejar la alimentación conectada permanentemente.

    \item \textbf{Varistor (VDR o MOV):} Es una resistencia dependiente de la tensión, ubicada en paralelo a la entrada. Funciona como protección contra picos transitorios. A tensiones de trabajo normales presenta una resistencia casi infinita (interruptor abierto), pero ante un pico o sobretensión extrema su resistencia cae drásticamente. Esto provoca un cortocircuito intencionado que absorbe el pico y obliga a fundir el fusible, salvando a los transistores y otras partes delicadas del equipo.

    \item \textbf{Termistor NTC:} Es un componente con un coeficiente de temperatura negativo diseñado para frenar el intenso pico de corriente provocado por los condensadores al encender el equipo. Al momento de la conexión la NTC se encuentra fría y ofrece una resistencia elevada limitando el paso de corriente. Inmediatamente después, el propio paso de corriente la calienta, haciendo que su resistencia caiga a un valor casi despreciable y permitiendo el flujo eléctrico normal.

    \item \textbf{Rectificador en puente (Rectificador primario):} Transforma la tensión de corriente alterna de entrada en corriente continua pulsante. Está formado por cuatro diodos que pueden venir en componentes discretos individuales o en un solo encapsulado o pastilla integrada. Al funcionar en conmutación (abriendo y cerrando el paso a la corriente), los diodos generan ruidos y oscilaciones de alta frecuencia, por lo que frecuentemente se coloca una \textbf{red RC (resistencia y condensador en serie)} en paralelo a la salida del rectificador. Esta red ayuda a amortiguar la conmutación y a disipar los armónicos nocivos en forma de calor.
\end{itemize}

\newpage
\section*{Esquema Eléctrico del Circuito de Entrada}

A continuación, se presenta el diagrama esquemático que ilustra la disposición típica de estos componentes en la etapa de entrada de una SMPS. 

\begin{figure}[H]
    \centering
    \resizebox{\textwidth}{!}{
    \begin{circuitikz}[american, thick, scale=1.0, transform shape]
        
    % Entrada de red (Línea y Neutro)
    \draw (0,4) node[anchor=east] {\textbf{L}} to[short, o-] (0.5,4)
          to[fuse, l=Fusible] (2.5,4) node[circ]{} node[above]{a} -- (3,4);
    \draw (0,0) node[anchor=east] {\textbf{N}} to[short, o-] (3,0);

    % Varistor (VDR) para protección contra sobretensiones
    \draw (2.5,4) to[varistor, l=VDR] (2.5,0);

    % Termistor NTC para limitar la corriente inicial
    \draw (3,4) to[thermistor, l=NTC] (5,4);
    \draw (3,0) to[short] (5,0);

    % Resistencias de descarga
    \draw (5,4) to[short] (6,4);
    \draw (5,0) to[short] (6,0);
    \draw (5.5,4) node[circ]{} node[above]{b} to[R, l=$R_{desc}$] (5.5,0);

    % Filtro EMC - Condensador de seguridad Tipo X
    \draw (6,4) -- (7.5,4) to[C, l=$C_X$] (7.5,0);
    \draw (6,0) -- (7.5,0);

    % Filtro EMC - Bobinas acopladas en modo común
    \draw (7.5,4) -- (8,4) to[L, l=$L_1$] (10,4);
    \draw (7.5,0) -- (8,0) to[L, l=$L_2$] (10,0);
    % Líneas punteadas para indicar el acoplamiento magnético (núcleo)
    \draw[dashed] (8.8,2.1) -- (9.8,2.1);
    \draw[dashed] (8.8,1.9) -- (9.8,1.9);

    % Filtro EMC - Condensadores de seguridad Tipo Y a tierra
    \draw (10,4) to[short] (10.5,4) coordinate(c) node[circ]{} node[above]{c} to[C, l_=$C_{Y1}$] (10.5,2);
    \draw (10,0) to[short] (10.5,0) coordinate(h) node[circ]{} node[below]{d} to[C, l=$C_{Y2}$] (10.5,2);
    \draw (10.5,2) node[circ]{} -- (10.8, 2) node[ground]{}; % Punto de unión y tierra a la derecha

    % Puente rectificador de onda completa (dibujado manualmente con 4 diodos)
    \coordinate (br-left) at (12.0,2);
    \coordinate (br-right) at (15.0,2);
    \coordinate (br-top) at (13.5,3.5);
    \coordinate (br-bottom) at (13.5,0.5);

    % Etiquetas para los nodos del puente
    \node[above left] at (br-left) {h};
    \node[above right] at (br-top) {e};
    \node[right] at (br-right) {f};
    \node[below right] at (br-bottom) {g};

    % Dibujo del puente con los diodos escalados a un tamaño menor
    \begin{scope}
        \ctikzset{diodes/scale=0.6}
        \draw (br-left) to[D, l=$D_1$, *-*] (br-top);
        \draw (br-top) to[D, l=$D_2$, *-*] (br-right);
        \draw (br-left) to[D, l_=$D_3$, *-*] (br-bottom);
        \draw (br-bottom) to[D, l_=$D_4$, *-*] (br-right);
    \end{scope}

    % Conexiones AC desde el filtro al puente rectificador
    \draw (c) -| (br-top);
    \draw (h) -| (br-bottom);

    % Resistencia de carga entre las salidas de CC (nodos h y f)
    \draw (br-right) |- (16, 5);
    \draw (br-left) |- (16, -1);
    \draw (16, 5) to[R, l=$R_L$, i=$I_L$, *-*] (16, -1);

    \end{circuitikz}
    }
    \caption{Esquema integral del circuito de entrada, protección y rectificación de una Fuente Conmutada.}
\end{figure}

\section*{Análisis Teórico: El Efecto de Tierra Flotante}

Para comprender el comportamiento de las señales antes del puente rectificador a la frecuencia de la red eléctrica (típicamente 50\,Hz o 60\,Hz), supongamos que entre los terminales de entrada \textbf{L} y \textbf{N} se aplica una señal sinusoidal de amplitud $A$, es decir, $V_{in}(t) = A \sin(\omega t)$.

\subsubsection*{1. Comportamiento de los componentes a baja frecuencia}
A la frecuencia de la red eléctrica, el circuito se comporta de una manera muy específica:
\begin{itemize}
    \item \textbf{Bobinas acopladas ($L_1$ y $L_2$):} Su reactancia inductiva ($X_L = 2\pi f L$) a 50/60\,Hz es extremadamente baja. Para efectos prácticos de la señal de red, actúan como un cortocircuito.
    \item \textbf{Termistor NTC y Fusible:} Presentan una resistencia muy baja.
\end{itemize}
En consecuencia, la señal de voltaje que entra entre \textbf{L} y \textbf{N} llega prácticamente intacta a los nodos \textbf{c} y \textbf{d}. Es decir, la diferencia de potencial entre ellos es $V_{cd} \approx A \sin(\omega t)$.

\subsubsection*{2. El divisor de tensión capacitivo ($C_{Y1}$ y $C_{Y2}$)}
El punto clave del análisis radica en el nodo de \textbf{tierra}, el cual está conectado exactamente en el punto medio de los condensadores de seguridad tipo Y ($C_{Y1}$ y $C_{Y2}$). Dado que en un filtro EMC bien diseñado $C_{Y1} = C_{Y2}$, estos actúan como un divisor de tensión capacitivo perfectamente balanceado. 

Lo que se mida respecto a esta ``tierra'' dependerá de cómo esté conectado el equipo a la instalación eléctrica:

\begin{itemize}
    \item \textbf{Escenario A: Equipo sin conexión a tierra (Tierra flotante).} Si medimos respecto al nodo central del divisor (que suele estar conectado al chasis metálico de la fuente), el divisor parte la amplitud de la señal de entrada exactamente por la mitad:
    \begin{itemize}
        \item Señal en el nodo \textbf{c} respecto a tierra: $V_c = \frac{A}{2} \sin(\omega t)$
        \item Señal en el nodo \textbf{d} respecto a tierra: $V_d = -\frac{A}{2} \sin(\omega t)$
    \end{itemize}
    \textit{Nota de ingeniería:} Esta es la razón por la cual, si se toca el chasis metálico de un equipo electrónico con la conexión a tierra interrumpida, se siente un leve hormigueo o descarga. El chasis está flotando a la mitad del voltaje de la red debido a estos condensadores.

    \item \textbf{Escenario B: Equipo con conexión a tierra real.} En la mayoría de las instalaciones eléctricas, el cable Neutro (\textbf{N}) se une físicamente a la toma de Tierra en el cuadro principal. Si este es el caso:
    \begin{itemize}
        \item El nodo \textbf{N} y la tierra están al mismo potencial ($0\,\text{V}$).
        \item Como $L_2$ es casi un cortocircuito a baja frecuencia, la señal en el nodo \textbf{d} respecto a tierra será prácticamente nula: $V_d \approx 0\,\text{V}$.
        \item El nodo \textbf{L} lleva toda la oscilación. Por tanto, la señal en el nodo \textbf{c} respecto a tierra será la onda completa: $V_c \approx A \sin(\omega t)$.
    \end{itemize}
\end{itemize}

\subsubsection*{3. Etapa de Rectificación de Onda Completa}
El puente rectificador está formado por los diodos $D_1$, $D_2$, $D_3$ y $D_4$. Su funcionamiento durante un ciclo completo de la señal alterna que llega desde el filtro EMC es el siguiente:
\begin{itemize}
    \item \textbf{Semiciclo positivo:} Cuando el nodo \textbf{e} es positivo con respecto al nodo \textbf{g}, la corriente fluye a través del diodo $D_2$, pasa por la resistencia de carga $R_L$ en el sentido del nodo \textbf{f} al nodo \textbf{h}, y retorna al nodo \textbf{g} a través del diodo $D_3$.
    \item \textbf{Semiciclo negativo:} Cuando la polaridad se invierte (el nodo \textbf{g} se hace positivo con respecto al nodo \textbf{e}), la corriente fluye a través del diodo $D_4$, atraviesa la carga $R_L$ exactamente en el mismo sentido que en el caso anterior (de \textbf{f} a \textbf{h}), y retorna al nodo \textbf{e} a través del diodo $D_1$.
\end{itemize}
De esta manera, la corriente sobre la carga $R_L$ circula siempre en una sola dirección, logrando la rectificación de onda completa y estableciendo al nodo \textbf{f} como el terminal positivo ($+V_{DC}$) y al nodo \textbf{h} como el retorno negativo (GND o $-V_{DC}$).

\subsubsection*{4. Función de las inductancias $L_1$ y $L_2$ (Filtro de Modo Común)}
Las bobinas $L_1$ y $L_2$ conforman lo que se conoce como un \textbf{choque de modo común} (Common-Mode Choke). Su función principal es suprimir el ruido electromagnético de alta frecuencia generado por la conmutación de la propia fuente, evitando que este contamine la red eléctrica externa, mientras permiten el paso libre de la energía a la frecuencia de red (50/60\,Hz). Su funcionamiento se basa en la dirección de la corriente y el acoplamiento magnético en el núcleo:
\begin{itemize}
    \item \textbf{Señales en Modo Diferencial (Energía de red):} La corriente alterna normal entra por la línea (\textbf{L}), atraviesa $L_1$, alimenta el circuito, y retorna por el neutro (\textbf{N}) atravesando $L_2$ en sentido opuesto. Al estar enrolladas estratégicamente sobre el mismo núcleo de ferrita, los campos magnéticos generados por estas dos corrientes opuestas de baja frecuencia se anulan entre sí. Como resultado, la impedancia magnética neta es casi cero (las bobinas actúan como simples cables) y el núcleo no se satura, sin importar cuánta corriente consuma la fuente.
    \item \textbf{Señales en Modo Común (Ruido electromagnético):} El ruido de alta frecuencia (típicamente entre 10\,kHz y varios MHz), originado por los veloces cambios de encendido y apagado del transistor principal de la SMPS, tiende a propagarse en la misma dirección simultáneamente por ambos cables (\textbf{L} y \textbf{N}) buscando retornar a través de tierra o radiar al ambiente. Para estas corrientes que fluyen en el mismo sentido, los campos magnéticos en $L_1$ y $L_2$ se suman constructivamente. Esto dota a las bobinas de una reactancia inductiva extremadamente alta, actuando como un "muro" o barrera que bloquea de forma efectiva la propagación de estas interferencias.
\end{itemize}

\subsubsection*{5. ¿Forman $R_{desc}$ y $C_X$ un filtro pasa bajo?}
Es una intuición común pensar que la resistencia de descarga ($R_{desc}$) y el condensador $C_X$ conforman un filtro RC pasa bajo. Sin embargo, \textbf{esto es un error conceptual}. 

Para que exista un filtro RC pasa bajo que atenúe el ruido de la línea, la resistencia tendría que estar conectada en \textit{serie} con el flujo de corriente principal. Colocar una resistencia en serie en la entrada de una fuente de alimentación es inviable, ya que toda la corriente de consumo pasaría por ella, disipando una enorme cantidad de energía en forma de calor ($P = I^2R$) y reduciendo drásticamente la eficiencia.

En la realidad, el filtrado EMI de una fuente conmutada no es de tipo RC, sino de tipo \textbf{LC (Inductor-Condensador)}. La atenuación real de alta frecuencia se logra mediante la combinación de las bobinas $L_1$ y $L_2$ (que se oponen a los cambios rápidos de corriente) y los condensadores $C_X$ y $C_{Y}$ (que cortocircuitan el ruido restante).

\textbf{Entonces, ¿cuál es la función de $R_{desc}$?} \\
Dado que $R_{desc}$ está conectada en \textit{paralelo} con $C_X$ y tiene un valor óhmico muy elevado (típicamente del orden de los Megaohmios, $\text{M}\Omega$), su impacto durante el funcionamiento normal es prácticamente nulo y no interviene en el filtrado. Su único y vital propósito es la \textbf{seguridad humana}. 

Cuando el equipo se desconecta del tomacorriente, el condensador $C_X$ queda cargado con la tensión de red. Si no existiera $R_{desc}$, esa carga eléctrica podría permanecer almacenada por horas. Si un usuario tocara las clavijas del enchufe desconectado, recibiría una fuerte descarga. $R_{desc}$ proporciona un camino para que el condensador se descargue de forma segura en una fracción de segundo tras la desconexión de la red.

\subsubsection*{6. Dispositivos de protección: Varistor (VDR) y Termistor NTC}
Además del filtrado, la etapa de entrada cuenta con componentes críticos para proteger la circuitería interna de la fuente frente a anomalías en la red y transitorios mecánicos de encendido:
\begin{itemize}
    \item \textbf{Varistor (VDR o MOV):} Actúa como la primera línea de defensa contra picos de sobretensión transitorios (como los inducidos por descargas atmosféricas o anomalías en la red). Se conecta en paralelo a la entrada. En condiciones normales, su resistencia es casi infinita. Sin embargo, si la tensión supera su umbral de diseño, su resistencia cae abruptamente, absorbiendo (cortocircuitando) la energía del pico transitorio. Si la sobretensión es muy alta o sostenida, este cortocircuito intencional provocará que el fusible principal se funda, desconectando y salvando de forma permanente al resto del equipo.
    \item \textbf{Termistor NTC (Coeficiente de Temperatura Negativo):} Su propósito es proteger a los diodos del puente rectificador frente al enorme pico de corriente inicial de encendido (\textit{inrush current}). Al conectar el equipo a la red, los condensadores principales de filtro (descargados) se comportan por una fracción de segundo como un cortocircuito. El NTC se encuentra frío en el momento de la conexión y presenta una resistencia moderada (por ejemplo, entre $5\,\Omega$ y $10\,\Omega$), lo que frena este pico de corriente. Inmediatamente después, el paso continuo de la corriente de trabajo calienta el termistor, haciendo que su resistencia caiga a un valor cercano a cero ($<1\,\Omega$), permitiendo el flujo normal de energía sin causar caídas de tensión importantes ni penalizar la eficiencia de la fuente.
\end{itemize}

\end{document}